\section{Arbitrary barriers and the distinction between primal and dual movement}
\label{sec:pd-movement}

The movement bounds of Section~\ref{sec:exposed-movement} depend on
certificate ranks and use a specified symmetric-cone metric. In the
primal--dual product metric, a lower bound depending
only on the barrier parameter holds for every logarithmically
homogeneous barrier. We state this classical argument and combine it
with the earlier formulation bounds. The dimension bound and the
objective-weight calculations below also appear in the companion
manuscript~\cite[Sections~5, 9, and~10]{CentralPathCompanion}; we
include them to make precise the metric and target set of each bound.

\subsection{The feasible gap set in the product metric}
Let $F$ be a $\nu$-logarithmically homogeneous standard self-concordant
barrier on a proper cone $K$. Define the conjugate barrier $F_*$ (with
pairing $-\ip{s}{x}$) and the product barrier $\widehat F$ by
\[
 F_*(s)=\sup_{x\in\intr K}\{-\ip{s}{x}-F(x)\},\qquad
 \widehat F(x,s)=F(x)+F_*(s).
\]
Consider the strictly feasible conic pair
\[
 \min\{\ip{c}{x}:Ax=b,\ x\in K\},\qquad
 \max\{\ip{b}{y}:A^Ty+s=c,\ s\in K^*\}.
\]
Write $z(\eta)=(x(\eta),s(\eta))$ for its exact centers, with
$s(\eta)=-\eta^{-1}F'(x(\eta))$ and gap $\nu/\eta$.
Distances $d_{\widehat F}$ allow all curves in $\intr K\times\intr K^*$;
requiring them to satisfy the affine equations can only increase
distances. A dot denotes differentiation in $\lambda=\log\eta$.

\begin{proposition}[Primal and dual parts of the central velocity]
\label{prop:pd-velocity}
Put $H=F''(x)$, $g=H^{-1/2}F'(x)$, and let $\Pi$ project orthogonally
onto $H^{1/2}\ker A$. Then
\begin{align}
 \nu_{\rm P}(\eta)&:=\norm{\dot x}_{F,x}^2=\norm{\Pi g}^2,
 &\nu_{\rm D}(\eta)&:=\norm{\dot s}_{F_*,s}^2
      =\norm{(I-\Pi)g}^2,\notag\\
 \nu_{\rm P}(\eta)+\nu_{\rm D}(\eta)&=\nu.
 \label{eq:pd-velocity}
\end{align}
After removing redundant equality rows,
\[
 \nu_{\rm P}=\eta^2c^TMc,\qquad
 M=H^{-1}-H^{-1}A^T(AH^{-1}A^T)^{-1}AH^{-1}.
\]
For $p(\eta)=\ip{c}{x(\eta)}$ and $d(\eta)=\ip{b}{y(\eta)}$,
\begin{equation}\label{eq:pd-objective-progress}
 -\eta^2p'(\eta)=\nu_{\rm P}(\eta),\qquad
 \eta^2d'(\eta)=\nu_{\rm D}(\eta).
\end{equation}
If both objectives converge to the finite optimum $p^*$, then
\[
 p(\eta)-p^*=\int_\eta^\infty\frac{\nu_{\rm P}(u)}{u^2}\,du,
 \qquad
 p^*-d(\eta)=\int_\eta^\infty\frac{\nu_{\rm D}(u)}{u^2}\,du.
\]
\end{proposition}
\begin{proof}
Homogeneity and conjugacy give $Hx=-F'(x)$, $\norm g^2=\nu$, and
$F_*''(s)=\eta^2H^{-1}$. Differentiation gives
$\eta\dot s=F'(x)-H\dot x$.
Thus $v=H^{1/2}\dot x$ and $d_0=\eta H^{-1/2}\dot s$ add to $g$.
Primal and dual feasibility put these vectors in $H^{1/2}\ker A$ and
its orthogonal complement, respectively. This proves the speed formulas.
Differentiating primal stationarity gives $\dot x=-\eta Mc$;
$MHM=M$ proves the displayed Schur formula and the first identity in
\eqref{eq:pd-objective-progress}. Differentiate $p-d=\nu/\eta$ for the
second, and integrate under the stated convergence assumption.
\end{proof}
This classical orthogonal decomposition underlies the
constant-speed primal--dual path of Nesterov and
Todd~\cite[Theorems~4.1 and~5.2]{NT2002}.

\begin{theorem}[Classical distance to a feasible gap sublevel]
\label{thm:pd-gap-movement}
Let $z_0=z(\eta_0)$ have gap $\Delta_0=\nu/\eta_0$, and let
$\mathcal Z_\epsilon$ be the strictly primal--dual feasible pairs with
gap at most $0<\epsilon<\Delta_0$. Then
\begin{equation}\label{eq:pd-gap-distance}
 \sqrt{\nu/2}\log(\Delta_0/\epsilon)
 \leq d_{\widehat F}(z_0,\mathcal Z_\epsilon)
 \leq\sqrt\nu\log(\Delta_0/\epsilon).
\end{equation}
Consequently, let $m\geq1$ be an integer and $0<R<1$. If a
trajectory from $z_0$ to $\mathcal Z_\epsilon$ consists of $T$ rounds,
each of at most $m$ chords from $z$ to $z+h$ with
$\norm h_z\leq R$, where $\norm h_z^2=\widehat F''(z)[h,h]$, then
\begin{equation}\label{eq:pd-round-bound}
 T\geq\frac{\sqrt{\nu/2}\log(\Delta_0/\epsilon)}
                  {m\log(1/(1-R))}.
\end{equation}
Intermediate chord endpoints need not satisfy the affine equations.
\end{theorem}
\begin{proof}
At $z_c=z(\eta)$, conjugacy gives $F_*'(s_c)=-\eta x_c$ and
$\widehat F(z_c)=\nu\log\eta-\nu$. For any feasible $z=(x,s)$,
$\ip{x-x_c}{s-s_c}=0$, since the first difference belongs to
$\ker A$ and the second to $\operatorname{range}A^T$. Hence
\[
 D\widehat F(z_c)[z-z_c]
 =-\eta\{\ip{s_c}{x-x_c}+\ip{x_c}{s-s_c}\}
 =-\eta\{\ip{x}{s}-\nu/\eta\}.
\]
At $\eta=\nu/\epsilon$, convexity shows that $z_c$ minimizes
$\widehat F$ on $\mathcal Z_\epsilon$. Every target therefore has
barrier increase at least $\nu\log(\Delta_0/\epsilon)$ from $z_0$.
The ambient dual norm of $D\widehat F$ is exactly $\sqrt{2\nu}$;
integration along any joining curve proves the lower bound. The central
path has constant speed $\sqrt\nu$ by
Proposition~\ref{prop:pd-velocity}, proving the upper bound.
For a chord $z+t h$ with $\norm h_z=r<1$, self-concordant Hessian
comparison bounds its speed by $r/(1-tr)$; integration gives length
at most $-\log(1-r)$. Sum over the chords.
\end{proof}
This argument is Theorem~5.1(c) of Nesterov and Todd; the chord
comparison and the primal--dual geometry are also theirs
\cite[Section~3 and Theorems~5.1--5.2]{NT2002}.
By the triangle inequality, an initial metric error $D_0$ subtracts
$D_0$ from the numerator's distance bound. The theorem says nothing
about unrestricted long moves, primal-only target sets, or query counts.

\subsection{A dimension bound for arbitrary cone dictionaries}
\begin{theorem}[Barrier parameter bound from dimensions alone]
\label{thm:dimension-pd-charge}
Let a compact full-dimensional body $C\subset\R^D$, $D\geq1$, have
an exact affine lift over a product of proper cones. After minimal-face
reduction (Lemma~\ref{lem:reduction}) and deletion of zero factors, suppose all factors have
dimension at most an integer $d\geq2$. Define
\[
 \Psi_d(h)=2\lfloor h/d\rfloor+\min\{h\bmod d,2\}.
\]
Every possibly coupled logarithmically homogeneous standard
self-concordant barrier on this reduced product has
\begin{equation}\label{eq:dimension-pd-charge}
 \nu\geq\Psi_d(D+1)\geq 2(D+1)/d.
\end{equation}
Thus the coefficient $\sqrt{\nu/2}$ in
\eqref{eq:pd-round-bound} is at least $\sqrt{\Psi_d(D+1)/2}$.
\end{theorem}
\begin{proof}
The total dimension $M$ of the reduced product is at least $D+1$, by
Theorem~\ref{thm:minimal-dimension-rigidity}. This also follows from
the ordinary rank $D+1$ of the body's affine slack operator after
the reductions of Lemma~\ref{lem:reduction}.
Let $r_0$ factors be rays and $q_0$ have dimension at least two.
In each nonray factor a two-dimensional linear section through an
interior point is a proper planar cone, hence an orthant after a linear
change of coordinates. Their product with all the rays is an interior
orthant section of dimension $r_0+2q_0$. Restriction of the barrier
preserves its homogeneity degree and self-concordance, so the sharp
orthant lower bound gives $\nu\geq r_0+2q_0$
\cite{NN1994,Hildebrand2013}. The argument permits arbitrary coupling.
The cap $d$ gives $r_0+dq_0\geq M\geq D+1$.
For $h=ad+b$, $0\leq b<d$, the minimum of $r+2q$ over nonnegative
integers with $r+dq\geq h$ is $2a+\min(b,2)$: for $q\leq a$ its
minimum is $h-(d-2)q$, while $q\geq a+1$ costs at least $2a+2$.
Finally $\min(b,2)\geq2b/d$.
\end{proof}
The function $\Psi_d$ is the exact integer minimum implied by these
two bounds; we do not claim
that some lift attains it for every body. No definability or contact
regularity is needed. For a product of balls it
uses $D=\sum_a s_a$, whereas on a positively curved
$D$-body with definable factors, Theorem~\ref{thm:curvature} gives the
usually stronger bound
\[
 \nu\geq2\left\lceil\frac{D-1}{d-2}\right\rceil
 \qquad(d\geq3).
\]
Under the hypotheses of Theorem~\ref{thm:smooth-lorentz}, the exact
minimum number $\sum_a h_a$ of Lorentz factors instead gives
$\nu\geq2\sum_a h_a$ and coefficient $\sqrt{\sum_a h_a}$.
For the single shared product-ball cone below, the optimal parameter is
$k+1$ and the coefficient $\sqrt{(k+1)/2}$.
All these bounds concern barriers on the ambient cone and its
conjugate, not the restricted standard barrier on a fixed-scale slice.

\subsection{Primal and dual speeds on shared product balls}
For positive dimensions $s_a$, consider
\[
 \mathcal H=\{(t,y_1,\ldots,y_k):t\geq\norm{y_a}_2\ (1\leq a\leq k)\}.
\]
Theorem~\ref{thm:shared-spectral-parameter}, specialized to vector rows,
gives the optimal $(k+1)$-logarithmically homogeneous barrier
\begin{equation}\label{eq:pd-product-barrier}
 F(t,y)=-\sum_a\log(t^2-\norm{y_a}^2)+(k-1)\log t.
\end{equation}
Its fixed-scale restriction is $f(y)=-\sum_a\log(1-\norm{y_a}^2)$.
For unit vectors $u_a$, weights $w_a\geq0$, and objective
$-\sum_a w_a\ip{u_a}{y_a}$ on $t=1$, the centers are
\begin{equation}\label{eq:pd-product-centers}
 y_a(\eta)=r(\eta w_a)u_a,\qquad
 r(z)=\frac{z}{\sqrt{1+z^2}+1}.
\end{equation}
Indeed $2r/(1-r^2)=\eta w_a$. Differentiation in $\log\eta$ gives
$\dot r=r(1-r^2)/(1+r^2)$; the radial Hessian is
$2(1+r^2)/(1-r^2)^2$. Consequently
\begin{equation}\label{eq:pd-product-effective}
 \nu_{\rm P}(\eta)=q_{\rm eff}(\eta)
 :=\sum_a\left(1-\frac1{\sqrt{1+(\eta w_a)^2}}\right),\qquad
 \nu_{\rm D}(\eta)=k+1-q_{\rm eff}(\eta).
\end{equation}
Thus $q_{\rm eff}$ smoothly counts the factors whose weight is
significant at scale $1/\eta$; we call $w_a$ visible at $\eta$ when $\eta w_a\geq1$.
Let $L_{\rm P}(\eta_0,\eta_1)$ and $L_{\rm D}(\eta_0,\eta_1)$ be the
primal and dual central-path lengths between $\eta_0$ and $\eta_1$, in
the metrics of $F$ and $F_*$.
With $w_1=1$ and all other weights zero, any interval
$0<\eta_0<\eta_1$ satisfies
\[
 L_{\rm P}\leq\log(\eta_1/\eta_0),\qquad
 L_{\rm D}\geq\sqrt{k}\log(\eta_1/\eta_0),\qquad
 d_{\widehat F}(z(\eta_0),z(\eta_1))
 \geq\sqrt{(k+1)/2}\log(\eta_1/\eta_0).
\]
From the primal analytic center $\eta=0$ to $\eta_1\geq1$, the primal
length is at most $1/\sqrt2+\log\eta_1$, since
$q_{\rm eff}(\eta)\leq\eta^2/2$ near zero.
A variant has a unique optimizer. For fixed
$0<\epsilon\leq1$, set $\nu=k+1$, $w_1=1$,
$w_a=\epsilon^2/\nu^2$ for $a>1$, and $\eta_1=\nu/\epsilon$.
Then all weights are positive, the final gap $\nu/\eta_1$ is
$\epsilon$, and throughout $1\leq\eta\leq\eta_1$,
\[
 q_{\rm eff}(\eta)\leq1+\frac{\epsilon^2}{2\nu},\qquad
 L_{\rm P}(1,\eta_1)
 \leq\sqrt{1+\frac{\epsilon^2}{2\nu}}\log(\nu/\epsilon).
\]
Thus, even with a unique optimizer, the barrier parameter alone gives
no lower bound on primal distance uniform over objectives. The weights here depend on $\epsilon$; for any fixed
instance with all $w_a>0$, $q_{\rm eff}(\eta)\to k$ as $\eta\to\infty$.
The general distinction between primal and product geometry is classical
\cite{NT2002,NesterovNemirovski2008}; the explicit specializations above
are also recorded in~\cite[Section~9]{CentralPathCompanion}.

\subsection{Accuracy and movement from the objective-weight distribution}
Conversely, the objective-weight distribution determines primal
accuracy and primal central length up to constant factors. Write
\[
 Q_j(\eta)=\sum_a\min\{(\eta w_a)^j,1\}\quad(j=1,2),\qquad
 E(\eta)=\sum_a w_a[1-r(\eta w_a)].
\]
Here $E(\eta)$ is the primal objective error of the center. For $z\geq0$, direct scalar calculation gives
\begin{align*}
 (1-1/\sqrt2)\min(z^2,1)
 &\leq1-(1+z^2)^{-1/2}\leq\min(z^2/2,1),\\
 (2-\sqrt2)\min(z,1)
 &\leq z[1-r(z)]\leq\min(z,1).
\end{align*}
For the first lower bound, use
$[1-(1+z^2)^{-1/2}]/z^2=
[\sqrt{1+z^2}(1+\sqrt{1+z^2})]^{-1}$ on $(0,1]$, and
monotonicity on $[1,\infty)$. For the second, use monotonicity of
$1-r(z)$ below one and of $1+z-\sqrt{1+z^2}$ above one.
Both lower constants are attained at $z=1$. Summation and integration
give
\begin{align}
 (2-\sqrt2)Q_1(\eta)/\eta&\leq E(\eta)\leq Q_1(\eta)/\eta,
 \label{eq:pd-distribution-error}\\
 \sqrt{1-1/\sqrt2}\,\mathcal J(\eta_0,\eta_1)
 &\leq L_{\rm P}(\eta_0,\eta_1)\leq\mathcal J(\eta_0,\eta_1),
 \quad
 \mathcal J=\int_{\eta_0}^{\eta_1}\sqrt{Q_2(\eta)}\frac{d\eta}{\eta}.
 \label{eq:pd-distribution-length}
\end{align}
The integral converges at $\eta_0=0$. After sorting the positive
weights, on an interval with $j$ weights satisfying $\eta w_a\geq1$,
write $B_j^2=\sum_{a>j}w_a^2$. Then $Q_2=j+\eta^2B_j^2$ and,
when $j,B_j>0$, an antiderivative is
\[
 \sqrt{j+\eta^2B_j^2}+\sqrt j\log
 \frac{\eta B_j}{\sqrt j+\sqrt{j+\eta^2B_j^2}}.
\]
The cases $j=0$ and $B_j=0$ give $\eta B_0$ and
$\sqrt j\log\eta$, respectively.

\begin{proposition}[Head and tail estimates]
\label{prop:pd-effective-support}
Assume $1\geq w_1\geq\cdots\geq w_k\geq0$,
$0<\epsilon\leq1$, and $1\leq m\leq k$.
If $\sum_{a>m}w_a\leq\epsilon/2$, the center at
$\eta_f=2m/\epsilon$ is primal $\epsilon$-accurate and
\begin{equation}\label{eq:pd-head-tail}
 L_{\rm P}(0,\eta_f)
 \leq\sqrt{(m+\epsilon^2/4)/2}
       +\sqrt{2m}\log(2m/\epsilon).
\end{equation}
Alternatively, if $\nu=k+1$ and
$\sum_{a>m}w_a^2\leq\epsilon^2m/\nu^2$, then at
$\eta_g=\nu/\epsilon$ the central primal--dual gap is $\epsilon$ and
\begin{equation}\label{eq:pd-full-gap-tail}
 L_{\rm P}(0,\eta_g)
 \leq\sqrt{(m+m\epsilon^2/\nu^2)/2}
       +\sqrt{3m/2}\log(\nu/\epsilon).
\end{equation}
\end{proposition}
\begin{proof}
The head error is at most $m/\eta_f$, the tail error at most its
weight sum. For $\eta\leq1$, speed is at most
$\eta\norm w_2/\sqrt2$, contributing $\norm w_2/\sqrt2$ to
length. For $1\leq\eta\leq\eta_f$ use
$q_{\rm eff}\leq Q_1\leq m+\eta\sum_{a>m}w_a\leq2m$,
and $\norm w_2^2\leq m+\epsilon^2/4$.
For the second claim, throughout $0<\eta\leq\eta_g$,
$q_{\rm eff}\leq m+(\eta^2/2)\sum_{a>m}w_a^2\leq3m/2$;
the same early-time integration and
$\norm w_2^2\leq m+m\epsilon^2/\nu^2$ prove the result.
\end{proof}
The second estimate bounds only the primal part of the central path to
gap $\epsilon$; by \eqref{eq:pd-product-effective}, the remaining
squared speed $k+1-q_{\rm eff}$ is dual. All tail
weights can be positive in either estimate. If at least $m$ weights are
visible throughout an interval, \eqref{eq:pd-distribution-length} gives
the matching lower bound $\sqrt{(1-1/\sqrt2)m}\log(\eta_1/\eta_0)$ on
arc length.

These arc estimates also give an explicit bounded-move schedule. Set
$\mathcal L(\eta)=\int_0^\eta\sqrt{q_{\rm eff}(u)}\,du/u$ and
choose successive centers at increments of at most $\log(1+R)$ in
$\mathcal L$, where $0<R<1$. For a nonzero objective this scalar
function is strictly increasing, so the parameters are uniquely
specified. Consecutive centers are at distance at most the arc length
between them, so the standard self-concordant comparison
$\log(1+\norm{y-x}_x)\leq d_f(x,y)$~\cite[Section~3]{NT2002} bounds
each chord's starting Dikin norm by $R$.
The number of chords is at most
$\lceil L_{\rm P}/\log(1+R)\rceil$. Evaluating and numerically
inverting $\mathcal L$, reading the weights, and computing the centers
have separate costs; the schedule bounds movement, not running time.

As an example with separated scales, take $w_a=2^{-(a-1)}$ and
$m=\lceil\log_2(4/\epsilon)\rceil\leq k$. The first estimate gives
$L_{\rm P}(0,2m/\epsilon)=O(\log^{3/2}(1/\epsilon))$ as
$\epsilon\downarrow0$. This is also the order of this central arc:
on $2^{j-1}\leq\eta\leq2^j$, $1\leq j\leq m$,
\[
 j\leq Q_2(\eta)\leq j+4/3.
\]
Each interval contributes $\Theta(\sqrt j)$, giving
$\Theta(m^{3/2})$ up to $2^m$; the remaining interval to
$2m/\epsilon$ contributes $O(\sqrt m\log m)$ because
$q_{\rm eff}\leq2m$ there. This conclusion concerns central arc
length and requires $k\geq m$, so $k$ grows as $\epsilon\downarrow0$;
for fixed finite $k$, all positive weights eventually become visible.
These distribution estimates and their spectral versions are companion
results~\cite[Section~5]{CentralPathCompanion}. Here they show that
neither the ambient barrier parameter nor the factor count can replace
the primal movement of a given formulation.

\begin{remark}[A pointwise obstruction to a parameter-only argument]
\label{rem:pd-pointwise-obstruction}
Dikin-ellipsoid containment, determinant growth, and the gradient
inequality alone do not give a primal distance bound depending only on
the parameter. In the cube at
$x_\delta=(1-\delta)\mathbf1$, $0<\delta\leq1$, put
$P=\mathbf1\mathbf1^T/k$ and
\[
 H_\delta=\frac{2}{k\delta^2}P+\frac{2}{\delta^2}(I-P),\qquad
 g_\delta=\frac{\mathbf1}{k\delta}.
\]
Then $(H_\delta^{-1})_{ii}=\delta^2(1-1/(2k))$, so its unit
ellipsoid fits inside every cube facet, while
$\norm{\mathbf1}_{H_\delta}=\sqrt2/\delta$,
$\det H_\delta=2^k/(k\delta^{2k})$, and
$g_\delta^TH_\delta^{-1}g_\delta=1/2$.
We do not claim that these pointwise data come from a global barrier.
The product-ball examples above already rule out such a bound uniform
over objectives; these matrices show that the weaker pointwise tests
cannot restore it. Neither construction decides whether such a bound
holds for all barriers when objectives are uniformly nondegenerate.
\end{remark}
